\chapter{Introduction}
\label{ch:introduction}

Social networks are nowadays largely used by people, allowing users to
interact and discuss the most different topics, but at the same time in these
platforms the polarization between users is constantly growing and several
studies have been conducted about it
\cite{Garimella2018}\cite{Guerra2013}\cite{conover2011political}\cite{gruzd2014investigating},
some of which focused on the recent debate about COVID-19
\cite{Jiang2021}\cite{green2020elusive}\cite{jiang2020political}\cite{lang2021maskon}
and vaccination \cite{Cossard2020}.

Polarization is the social phenomenon according to which people tend do
separate in opposing communities with few people remaining neutral
\cite{Guerra2013}. A close phenomenon is that of the Echo Chambers, groups in
which people that have the same opininions enforce their respective ideas
\cite{Garimella2018}, a concept very similar to the definition of polarization
as given in \cite{sunstein1999law}: ``group polarization arises when members of
a deliberating group move toward a more extreme point in whatever direction is
indicated by the members' predeliberation  tendency. `[L]ike polarized
molecules, group members become even more aligned in the direction they were
already tending.'\cite{turner1987rediscovering}''

In this document we aim at finding a method for finding Echo Chambers, by
analyzing data retrieved from social medias like Twitter and Reddit: we define
the \acrfull{ECP} and the \acrfull{D-ECP} and propose techniques for solving
and approximating them.

Both these $2$ problems are defined on a signed graph, i.e. distinguishing
between \emph{friendly} and \emph{hostyle} interactions between the users.
Differently from previous studies of polarization on signed graph
\cite{xiao2020searching}, we define
and incorporate in our problems the ideas of \emph{contents} (the piece of
information which is discussed) and \emph{threads} (the local community
discussing a content).

\section{Background}
\label{sec:background}

A \emph{graph} is a collection of \emph{vertices} or \emph{nodes} $V$ and
\emph{edges} or \emph{links} $E$ between the nodes, representing relationships
between entities; this structure turns out to be very useful in
representing many interesting concepts from social sciences, biology, physics,
chemistry and geography (\autoref{fig:tex/img/sample-graph})\cite{Newman2018}\cite{Menczer2020}.

\begin{figure}
	\centering
	\includegraphics[width=0.6\linewidth]{tex/img/sample-graph.png}
	\caption[Retweet network during 2010 midterm elections]{The retweet network of posts regarding US during 2010 midterm
		elections. Red and blue nodes are associated with conservative and
		progressive users, respectively \cite{Menczer2020}}%
	\label{fig:tex/img/sample-graph}
\end{figure}

\paragraph{Different types of graphs}%
\label{par:different_types_of_graphs}

In its simplest form graph are \emph{undirected} and \emph{unweighted}. In an
\emph{undirected} graph relationships are bi-directional, while in directed graph
the order of nodes in a link reflects the
direction (i.e. an edge $e_{ij} $ is different from an edge $e_{ji} $). A weighted
graph instead associates a weight $\omega $ to each edge
\cite{Menczer2020}\cite{AlbertLaszloNortheasternUniversity2016}.

Sometimes also edges are allowed to be either positive or negative (for example when
definying the relationships in a acquaintance network): these king of networks
are usually called \emph{signed graphs} \cite{Newman2018}.

\bigskip

In the rest of the document we will abuse notation and refer to vertices both
as $v_{i} \in V $ and $i \in V$; similarly we will refer to edges both as
$e_{ij} \in E $ and as $ij \in E$.

\bigskip

Representing relationsips between entities which span over more than $2$
dimensions requires the definition of an ever more complex structure, the
\emph{multiplex graphs}, which is a collection of graphs (referred to as
\emph{layers}) over the same set of vertices, where each layer could represent
a different type of connection \cite{Newman2018}.

\section{Problem}
\label{sec:problem}

\subsection{The Interaction Graph}
\label{sub:interaction-graph}

The \emph{Interaction Graph} $G$ is the graph which encodes the informations
regarding the interactions between the users.

In this graph each users is associated to a vertex $v \in V$ and each edge to
an interaction between the $2$ corresponding users it links. For this reason we
will sometime refer to vertices as users in the rest of the document.

$G = (V, E^{+}, E^{-})  $ is a
signed and weighted graph, the weights being in the interval $[-1, +1]$,
corresponding to positive and negative interactions, meaning that for smaller
the value of the weight the interaction will more "negative".

The \emph{Interaction Graph} is also directed, so that an edge from vertex $v_{i}
$ to vertex $v_{j} $ corresponds to a reply from user $v_{i} $ to $v_{j} $.

Let a content $C$ be any kind of resource that triggers a discussion in one or
more threads $T$. The set of threads associated to $C$ is denoted as
$\mathcal{T}_{C} $. A content is usually represented by a newspaper article and
it is identified by its url, e.g.

	{\footnotesize
		\begin{center}
			\url{https://www.nytimes.com/2021/03/04/us/richard-barnett-pelosi-tantrum.html}
		\end{center}
	}

A corresponding thread may then be, for example, a user posting and commenting
the same url on its Twitter account (see \autoref{fig:twitter-thread}), thus
generating a discussion.

\begin{figure}
	\centering
	\includegraphics[width=0.6\linewidth]{tex/img/twitter_thread.png}
	\caption[Thread-content distinction example from Twitter]{An thread associated to the mentioned New York Times article}%
	\label{fig:twitter-thread}
\end{figure}

The \emph{Interaction Graph} is a \emph{multiplex graph}, each layer being
represented by a thread $T$ whose edges are the interactions happening in it.
Note also that for this reason each of the layers can contain more than one
edge between $2$ users, as each pair of users can reply to each other more than
once.

We will also use $\mathcal{C} $ for denoting the set of contents.

An example of \emph{Interaction Graph} can be seen in
\autoref{fig:interaction-graph-example}.

\begin{figure}
	\begin{center}
		\begin{subfigure}[b]{0.3\textwidth}
			\centering
			\tikzfig{tex/tikz/graph_thread1}
			\caption{$T_{1} \in \mathcal{T}_{C_{1}} $}
			\label{fig:tex/tikz/graph_thread1.tikz}
		\end{subfigure}
		\begin{subfigure}[b]{0.3\textwidth}
			\centering
			\tikzfig{tex/tikz/graph_thread2}
			\caption{$T_{2} \in \mathcal{T}_{C_{1}} $}
			\label{fig:tex/tikz/graph_thread2.tikz}
		\end{subfigure}
		\begin{subfigure}[b]{0.3\textwidth}
			\centering
			\tikzfig{tex/tikz/graph_thread3}
			\caption{$T_{3} \in \mathcal{T}_{C_{2}} $}
			\label{fig:tex/tikz/graph_thread3.tikz}
		\end{subfigure}
	\end{center}
	\caption[Example \emph{Interaction Graph}]{An example of \emph{Interaction Graph}, green and red edges
		representing positive and negative interactions, respectively. It
		contains $3$ threads and $2$ contents, the first $2$ layers each being
		associated to a thread of content $C_{1} $, the last layer to a thread
		of content $C_{2} $}
	\label{fig:interaction-graph-example}
\end{figure}

\subsection{The problem definition}%
\label{sub:the_problem_definition}

The main goal of the research is finding echo chambers in social medias, more
specifically on the \emph{Interaction Graph} as defined in
\autoref{sub:interaction-graph}.

Our definition is based on the idea that echo chambers can be identified by
looking at contents which is generally highly debated (we will call
this type of content \emph{controversial}) but which is discussed with little
or no animosity in small subgraphs. This subgraphs are the \emph{Echo
	Chambers}.

\bigskip

Given an \emph{Interaction Graph} $G = (V, E^{+}, E^{-})$ on some contents
$\mathcal{C} $ and threads, let $\eta(T)$ the number of negative edges over the
total number of edges in the layer associated to thread $T$. Similarly,
let $\eta(C)$ be the fraction of negative edges in all threads associated to
content $C$.

\begin{definition}[Controversial thread]
	Let $\alpha \in [0,1]$. A thread (or content) is \emph{controversial} if
	$\eta(T) > \alpha$ (or, similarly, $\eta(C) > \alpha $). Conversely, a
	thread (or content) is \emph{non-controversial} if $\eta(T) \leq \alpha$
	($\eta(C) \leq \alpha$).
\end{definition}

Intutively, \emph{controversial} threads contain many negative
interactions. We denote as $\hat{\mathcal{C} } \subseteq \mathcal{C} $ the
set of \emph{controversial} contents.

\medskip

\emph{Echo Chambers} correspond to \emph{non-controversial} smaller subgraphs
(i.e. with few negative edges) discussing a
\emph{controversial} content.

More formally, let $T[U]$ the subgraph induced in the layer associated to
thread $T$ by the vertices $U$; let $|T^{+} [U]|$ and $|T^{-} [U]|$ its number
of positive and negative edges, respectively.

We define $\mathcal{S}_C (U)$ as the set of \emph{non-controversial} threads
induced by $U$, for \textit{controversial} contents, i.e.
	{\small
		\begin{equation}
			\mathcal{S} _{C} (U) = \{ T[U] \; s.t. \; T[U] \; non \;
			controversial, T \in \mathcal{T} _{C}, C
			\in \hat{\mathcal{C}}, U \subseteq V\}
		\end{equation}
	}

Thus $\mathcal{S} _C (U)$ will contain threads which are \emph{globally} non
controversial but it is defined only for contents that are \emph{globally}
controversial.

\medskip

We know define the Echo Chamber Score of a set of vertices $U$.

\begin{definition}[Echo Chamber Score]
	Let $U \subseteq V$ be a subset of vertices. Its Echo Chamber Score is

	\begin{equation}
		\label{eq:echo-chamber-score}
		\xi(U) = \sum^{}_{\mathcal{C} \in \mathcal{\hat{C}}} \sum^{}_{T[U] \in
		\mathcal{S} _{C} (U)} (|T^{+} [U]| - |T ^{-} [U]|)
	\end{equation}
\end{definition}

We can now define the \acrfull{ECP}.

\begin{problem}[\acrfull{ECP}]
Given an \emph{Interaction Graph} $G$ and $\alpha \in [0, 1]$ find a set of vertices $U \subseteq
	V$ maximizing the Echo Chamber Score (\autoref{eq:echo-chamber-score}).
\end{problem}

We will denote with $\hat{U}$ the set of users maximizing
\autoref{eq:echo-chamber-score} and with $\xi(G)$ its corresponding score, i.e.

\begin{align*}
	\hat{U} & \coloneqq \argmax_{U \subseteq V} \xi(U) & \xi(G) & \coloneqq
	\xi(\hat{U})
\end{align*}

\subsection{The Densest Echo Chamber Problem}%
\label{sub:the_densest_echo_chamber_problem}

The \acrshort{ECP} doesn't take into account the number of users producing a
certain score; this means that the set $U$ may involve disconnected and sparse
subgraphs, depending on the structure of the graph $G$.

For this reason it is interesting also to study another variant of the
\acrshort{ECP}, the \acrlong{D-ECP}, which we now define.

\begin{definition}[Echo Chamber Score]
	Let $U \subseteq V$ be a subset of vertices. Its Densest Echo Chamber Score is

	\begin{equation}
		\label{eq:densest-echo-chamber-score}
		\xi(U) = \sum^{}_{\mathcal{C} \in \mathcal{\hat{C}}} \sum^{}_{T[U] \in
		\mathcal{S} _{C} (U)} \frac{(|T^{+} [U]| - |T ^{-} [U]|)}{|U|}
	\end{equation}
\end{definition}

Similarly to the \acrshort{ECP} we can now define the corresponding problem

\begin{problem}[\acrfull{D-ECP}]
Given an \emph{Interaction Graph} $G$ and $\alpha \in [0, 1]$ find a set of vertices $U \subseteq
	V$ maximizing the Densest Echo Chamber Score (\autoref{eq:densest-echo-chamber-score}).
\end{problem}

\section{Goals}
\label{sec:goals}

In this document we address the following research questions:

\begin{enumerate}
	\item Can we solve the \acrlong{ECP} and the \acrlong{D-ECP}?
	\item Are the \acrlong{ECP} and the \acrlong{D-ECP} capable of finding echo
	      chambers inside social medias?
\end{enumerate}

\section{Research metodology}%
\label{sec:research_metodology}

% todo

\section{Structure of the thesis}
\label{sec:structure-thesis}

% todo

\section{About the thesis}
\label{sec:about-thesis}

The code which was used to obtain the results presented in the rest of the
document is available at the following GitHub address

\begin{center}
	\url{https://github.com/morpheusthewhite/master-thesis}
\end{center}
